% 微网四问：统一价值函数框架下的购电与储能调度（全面重写版）
\PassOptionsToPackage{fontset=none}{ctex}
\documentclass[withoutpreface,bwprint]{cumcmthesis}
\setCJKmainfont[Path=/System/Library/Fonts/Supplemental/,FontIndex=6,BoldFont=Songti.ttc,BoldFeatures={FontIndex=1}]{Songti.ttc}
\setCJKsansfont[Path=/System/Library/Fonts/,FontIndex=1]{STHeiti Medium.ttc}
\setCJKmonofont[Path=/System/Library/Fonts/,FontIndex=1]{STHeiti Light.ttc}
\setCJKfamilyfont{zhsong}[Path=/System/Library/Fonts/Supplemental/,FontIndex=6]{Songti.ttc}
\setCJKfamilyfont{zhhei}[Path=/System/Library/Fonts/,FontIndex=1]{STHeiti Medium.ttc}
\setCJKfamilyfont{zhkai}[Path=/System/Library/Fonts/Supplemental/,FontIndex=6]{Songti.ttc}
\setCJKfamilyfont{zhfs}[Path=/System/Library/Fonts/Supplemental/,FontIndex=6]{Songti.ttc}
\providecommand{\songti}{\CJKfamily{zhsong}}
\providecommand{\heiti}{\CJKfamily{zhhei}}
\providecommand{\kaishu}{\CJKfamily{zhkai}}
\providecommand{\fangsong}{\CJKfamily{zhfs}}
\setlength{\emergencystretch}{2em}
\usepackage{cumcm2026}
\usepackage{url}
\usepackage[section]{placeins}
\usepackage{colortbl}% 网格表表头底纹
\newcommand{\tabheadrow}{\rowcolor{cumcmgray}}
\newtheorem{proposition}{命题}
\graphicspath{{../../plot/}}
% 缺图时显式画占位框（交稿前 U6 终检改为硬失败，绝不静默丢图）
\newcommand{\safeincludegraphics}[2][width=0.9\textwidth]{%
  \IfFileExists{../../plot/#2.pdf}{\includegraphics[#1]{#2}}{%
  \fbox{\parbox[c][4cm][c]{0.8\textwidth}{\centering 缺图占位：\texttt{\detokenize{#2}}}}}}
\newcommand{\QoneCost}{35,126.95}
\newcommand{\QoneEnergy}{59,482.70}
\newcommand{\QtwoCost}{1356.6396}
\newcommand{\QthreeCost}{1302.4597}
\newcommand{\QfourTwoCost}{1427.0607}
\newcommand{\QfourThreeCost}{1371.6972}
\title{信息状态与未来费用函数：微网购电与储能调度的统一模型}
\tihao{C}\baominghao{0000}\schoolname{XX大学}
\membera{ }\memberb{ }\memberc{ }\supervisor{ }
\yearinput{2026}\monthinput{9}\dayinput{12}
\begin{document}
\maketitle
\begin{abstract}
\setstretch{1.23}
随着分布式光伏和用户侧储能持续接入，微网运行逐步转向源、荷、储与外部电网的协同调度。实际运行中，光伏出力、用户负荷与外网电价在时间上存在错位，购电计划需在未来状态尚未完全实现时提前确定，而预报与价格信息又在执行过程中逐步揭示，使采购调整与储能跨期配置相互影响。针对这类具有提前决策、信息持续更新和库存跨期传递特征的调控问题，本文以信息状态刻画各决策时刻可获得的观测、预报与交易权限，以库存未来费用函数度量储能电量的跨期机会成本，构建随信息扩展逐层递进的购电--储能联合决策框架。

\textbf{针对问题一}，首先根据微网能量流动关系，以全天计划购电费用最小为目标，综合考虑能量平衡、储能容量、充放电功率及日初日末库存一致等约束，建立确定性线性规划模型，求得最优全天计划购电量为 59482.70\,kWh，购电费用为 \textbf{35126.95 元}，较无储能直购降低 \textbf{26.90\%}。其次，为刻画储能电量的跨时段价值，以当前储电量为状态定义剩余最低费用函数 $F_t(e)$，证明其在有效域上具有凸分段线性结构，且库存边际价值随储电量增加而递减。最后，利用该性质建立连续状态动态规划独立复算，与线性规划最优值仅差 $7.3\times10^{-12}$ 元，为后续随机条件下的库存决策提供价值基础。

\textbf{针对问题二}，在问题一库存价值刻画的基础上，进一步考虑负载与光伏实际出力的不确定性。首先，将负载预测分解为日电量水平和归一化日内形状，并将历史负载、光伏预测误差按完整日路径配对，构造 30 条经验场景以保留误差的日内相关结构。其次，利用样本平均近似确定每天 0:00 的计划购电量，并通过 48\,h 前瞻与次日价值函数处理跨日库存价值；执行阶段将确定性费用函数推广为场景平均未来费用函数 $\bar H_t(e)$，据其边际结构形成动态库存保留水平，并依据实时观测逐 10\,min 控制储能。结果表明，2 月 1 日至 12 月 31 日实际费用为 \textbf{1352.96 万元}，较完美信息离线下界高 10.6\%，较同一预测信息下无储能计划方案降低 28.0\%。

\textbf{针对问题三}，在问题二全天计划购电策略的基础上，引入 0:00、6:00、12:00 和 18:00 发布的光伏预报及日内购电调整机制。首先，在各预报发布时刻利用最新光伏信息更新剩余时段场景，并以真实当前库存为初值滚动优化尚未执行的购电量；其次，在目标函数中统一计入计划购电、调增购电、调减违约及紧急购电费用，每轮仅提交至下一预报时刻，使后续计划能够随新信息继续修正。全年实际费用为 \textbf{1298.96 万元}，较问题二减少 54.01 万元，紧急购电费用由 78.23 万元降至 19.16 万元。进一步比较四种预报开放方案，新增 6:00、12:00 和 18:00 预报及其调整机会分别节省 3.83、\textbf{23.74}、5.09 万元，其中 12:00 的边际收益最高。

\textbf{针对问题四}，进一步考虑外网电价的实时波动，将问题二、三的购电策略推广至随机电价环境。首先，根据电价的跨日水平变化和日内峰谷特征，将电价分解为日价格水平与归一化日内形状；其次，将电价预测残差与负载、光伏残差按照同一历史信息状态进行三通道配对，构造保留供需--价格相关性的联合场景，并分别沿用问题二的 0:00 冻结采购权限和问题三的日内滚动调整权限求解。48\,h 主方法下，两类策略全年实际费用分别为 \textbf{1422.99 万元}和 \textbf{1369.30 万元}；在 24\,h 同口径比较中，与完美价格信息参照的费用差分别为 0.62\% 和 0.68\%。进一步证明，当全部价格输入及相应费用系数按同一正常数缩放时，最优采购解集、库存保留水平和执行动作保持不变。

综上，四类策略均以真实运行数据结算的实际费用作为评价依据。进一步的灵敏度分析与组件比较表明，场景数达到 20 后全年费用趋于稳定，48\,h 跨日价值处理在各类策略中均能降低实际费用；模型参数仅依据评分期前信息确定，全年回测严格遵循信息时序，验证了模型结果的稳定性与策略的时序可实施性。
\keywords{微网储能调度\quad 样本平均近似\quad 动态规划\quad 滚动优化\quad 库存价值函数}
\end{abstract}

% 一、问题重述
\section{问题重述}
\subsection{问题背景}
非孤岛式微网由小区负载、光伏发电、储能设备和外部电网共同构成。光伏出力具有明显的时变性，储能设备可以在电价较低或系统存在富余电量时充电，并在高价或供能不足时释放电量，从而在满足负载需求的前提下降低外网购电费用。微网调控需要根据负载、光伏出力、储能状态以及可获得的电价和预报信息，合理确定各时段购电量与储能充放电量。

题设储能最大容量为 12000\,kWh，运行储电量限制在 1200--10800\,kWh，最大充放电功率为 5000\,kW，2025 年 1 月 1 日 0:00 储电量为 6000\,kWh，充放电效率均为 90\%。附件 1 给出典型日的电价、负载及光伏预测功率；附件 2 给出 2025 年全年负载与光伏实际功率；附件 3 给出每天 0:00、6:00、12:00、18:00 发布的未来 24\,h 整点光伏预报；附件 4 给出全年实时电价。随着可用信息和交易权限增加，题目依次要求建立四类购电与储能调度策略。

\subsection{问题要求}
\textbf{问题 1}\quad 每日负载和电价保持一致，并给定一天的光伏预测曲线。要求每天 0:00 一次性确定全天计划购电策略，在各时段微网供能不低于负载、储能满足容量与功率限制且 0:00 与 24:00 储电量相同的条件下，使全天计划购电费用尽可能低，并给出指定时段购电量、全天购电量与购电费用，以及各指定区间的储能充放电量。

\textbf{问题 2}\quad 电价仍采用附件 1 给定的日内曲线，负载和光伏出力随日期变化。每天 0:00 确定全天计划购电量，计划确定后不再调整；实际运行中若微网供能出现缺口，则按交易时刻电价的 5 倍向外网紧急购电。需要结合附件 2 构造逐日购电策略，计算指定日期的计划购电、充放电和紧急购电结果，并保存 2 月 1 日至 12 月 31 日的完整策略。

\textbf{问题 3}\quad 在问题 2 基础上，每天 0:00、6:00、12:00 和 18:00 均可获得未来 24\,h 整点光伏预报。0:00 预报用于制定全天计划，后续预报可用于调整尚未执行的购电量。调减部分按交易时刻电价的 50\% 支付违约费用，超过原计划的调增部分按 1.5 倍交易时刻电价购入，总费用还包括正常计划购电和紧急购电费用。需要建立能够随预报更新进行调整的购电策略，并定量判断引入其他时刻预报的实际价值。

\textbf{问题 4}\quad 将外网电价进一步扩展为实时波动电价。结合附件 2 的实际负载与光伏数据、附件 3 的光伏预报以及附件 4 的全年电价，在波动电价条件下分别按照问题 2 和问题 3 的交易权限重新建立购电策略并计算相应结果。

% 二、问题分析
\section{问题分析}
\subsection{总体分析：信息状态与库存机会成本}
四个问题具有相同的物理约束。每一时段均需满足外网购电、光伏供能和储能放电与负载、储能充电之间的能量平衡；储电量按照充放电动作递推，并受容量与功率上限约束。随着题目推进，变化主要体现在决策前可获得的信息逐步丰富，以及购电计划由 0:00 一次确定扩展为允许日内调整。

储能状态使各时段决策相互关联。当前释放的库存能够降低本时段购电需求，同时会减少未来高成本缺口时可利用的电量。因此，仅记录一条计划库存轨迹不足以支持随机环境下的实时执行，还需要刻画不同库存状态对应的未来最低费用。

记 $\mathcal I_t$ 为时刻 $t$ 的信息状态，其中包含截至当前已经获得的负载、光伏和价格观测，已经发布的预测信息，当前储电量以及当前允许使用的交易权限。在给定 $\mathcal I_t$ 的条件下，以库存 $e$ 为状态构造未来费用函数，由其边际变化衡量库存电量的跨时段机会成本。问题 1 得到完整信息下的确定性价值函数 $F_t(e)$；问题 2 将其扩展为预测场景下的平均未来费用 $\bar H_t(e)$；问题 3 进一步随预报发布时刻 $\tau$ 更新；问题 4 再纳入未来价格信息。由此，四问可以在同一能量调度结构下随信息状态逐步扩展。

\subsection{问题 1 的分析}
问题 1 给定一天内用于决策的电价、负载与光伏预测曲线，各时段购电量、充放电量和库存之间均为线性关系，目标函数为各时段电价与购电量之和，因此建立确定性线性规划，在能量平衡、充放电功率、库存边界及日初日末库存相等的约束下求得全天计划的最低购电费用。

线性规划能够给出一组全局最优调度结果，但后三问还要处理执行过程中库存偏离计划值后的重新决策。为此，以当前储电量 $e$ 为状态，定义从时段 $t$ 开始至计划结束的剩余最低费用 $F_t(e)$，并建立连续状态 Bellman 递推。由于单段动作成本具有凸分段线性结构，$F_t(e)$ 在有效域上保持凸分段线性，可用折线断点和斜率精确表示，无需对储电量进行网格离散。LP 承担确定全天计划调度和最优费用的作用，连续状态 DP 提供库存状态的未来价值表示；两种独立求解路径的最优值一致性又构成对优化核的交叉验证。

\subsection{问题 2 的分析}
问题 2 中负载与光伏出力在 0:00 制定计划时尚未完全实现，需要同时处理计划采购的不确定性和实际运行阶段的储能控制。预测层利用附件 2 的历史数据描述可预测结构：将负载拆解为日电量水平与归一化日内形状，分别预测后重构全天计划曲线；光伏则使用决策时点以前已经完成的历史日构造基准预测。

预测误差在一天内具有连续性，负载与光伏误差之间也可能存在共同变化，因此场景层以历史完整日的预测残差为基本单元，并保留同一历史日的负载--光伏残差配对关系。决策时刻选取最近 $M=30$ 条已经完整实现的历史残差路径，叠加到当前点预测后形成经验场景集。

0:00 计划量一经制定便在全天计划内冻结，具有"先确定采购、后观察随机实现"的两阶段结构。采购层采用样本平均近似\cite{saa-book,kleywegt}，在多个历史场景下共同优化第一阶段计划购电量；场景内第二阶段变量用于评价储能追索和潜在紧急购电。由于该代理问题允许场景路径完整可见，其最优值用于计划选择与机会成本估计，不直接作为真实执行动作。

真实执行只使用当前已经观测的信息。对固定计划和每条场景逆向计算未来费用，并取等权平均得到 $\bar H_t(e)$。当出现供能缺口时，将当段紧急购电边际成本与未来库存边际价值进行比较，得到动态库存保留水平 $R_t$：库存高于该水平的可用部分可以释放，剩余部分为未来潜在高成本缺口保留。为减弱 24:00 截断造成的终端效应，主方法将前瞻范围扩展到次日，并以次日自由购电的价值函数给午夜库存定价。全年实际费用最终由真实负载、真实光伏和真实结算规则重新计算。

\subsection{问题 3 的分析}
问题 3 同时增加日内预报信息和购电调整权限。0:00 制定的原计划作为全天基准；6:00、12:00 和 18:00 获得新的未来 24\,h 光伏预报后，剩余净负荷分布和库存机会成本需要随最新信息重新评估。

在每个预报发布时刻，以真实当前储电量作为新的初始状态，利用截至该时刻已经发布的光伏预报和历史误差场景重新优化尚未执行的购电量，并在目标函数中统一计入调增、调减、紧急购电和跨日库存价值。本轮只正式提交至下一次预报更新时间之前的购电决策，更远时段保留为前瞻计划，待下一次预报到来后重新计算，从而保持整个滚动过程的时序一致性。

为回答是否需要引入其他时刻预报，固定预测方法、场景构造、终端价值和结算规则，按照 $\{0\}$、$\{0,6\}$、$\{0,6,12\}$、$\{0,6,12,18\}$ 的顺序逐步开放日内更新时刻，并以全年真实账单比较相邻方案。由于新增一个发布时刻同时带来新的光伏信息和一次对应的购电调整机会，相邻账单差表示该更新时刻在既定开放顺序下的联合边际实际收益。现有结果表明，12:00 更新的收益最高，其原因与午间光伏预报误差明显收缩以及当天仍保留较长的剩余执行窗口共同相关。

\subsection{问题 4 的分析}
问题 4 中未来电价也需要在购电决策前估计，库存未来费用因此同时取决于未来净负荷与价格路径。根据附件 4 的全年价格结构，将每日电价分解为日价格水平和归一化日内形状，前者刻画跨日整体价格变化，后者描述日内相对稳定的峰谷结构。

联合场景构造时，将价格预测残差与负载、光伏残差按照同一历史信息状态组成三通道完整路径，避免独立抽样破坏变量间的经验相关关系。4-2 沿用问题 2 的交易权限，0:00 计划购电量保持冻结，执行阶段随已经观测的价格信息更新库存价值；4-3 沿用问题 3 的交易权限，在各发布时刻同步更新光伏预报、价格预测与剩余采购。

当全部历史与实际价格、紧急及调整交易价格以及续存费率同步乘以同一正常数时，各优化问题的目标函数整体同比缩放而可行域保持不变，因此最优采购解集、库存保留水平和逐段执行动作均保持不变，费用按相同比例变化。这一齐次性质说明，统一改变绝对价格水平不会改变策略形状，真正影响调度时机的是相对价格结构及其可预测性。

\subsection{总体研究框架}
全文在纵向上依次经过预测与场景构造、采购优化、因果执行和实际结算；在横向上随着问题推进，决策环境由确定信息扩展至供需随机性、日内信息更新和价格不确定性。四问共享同一储能物理约束和基本优化结构，库存未来费用函数随信息状态扩展形成：
\begin{equation}
F_t(e)\;\longrightarrow\;\bar H_t(e)\;\longrightarrow\;\bar H_{t\mid\tau}(e)\;\longrightarrow\;\bar H_{t\mid\tau}(e;\,p),
\label{eq:chain}
\end{equation}
其中，$F_t(e)$ 描述问题 1 完整信息条件下的确定性剩余费用；$\bar H_t(e)$ 用于问题 2 的场景平均库存价值；$\bar H_{t\mid\tau}(e)$ 随问题 3 各日内预报时刻更新；$\bar H_{t\mid\tau}(e;\,p)$ 进一步纳入问题 4 的未来价格路径。问题 2--4 采用连续 walk-forward 方式评价，优化模型中的场景代理目标用于制定计划和估计库存机会成本，最终结果统一按真实实现的负载、光伏、电价以及题目规定的交易费用重新结算。

\begin{figure}[!htbp]\centering
\safeincludegraphics[width=0.96\textwidth]{fig_f01_framework}
\caption{四问的预测、采购与执行框架}
\label{fig:framework}
\end{figure}

% 三、模型假设；四、符号说明
\section{模型假设}
\begin{enumerate}
\item \textbf{时间离散与段内恒定}。将附件中的功率序列按 10\,min 粒度离散，一天共 $T=144$ 个时段；每一时段内负载、光伏和储能充放电功率视为恒定，功率 $P$（kW）对应的时段电量为 $P/6$（kWh）。
\item \textbf{当前状态可观测}。执行进入时段 $t$ 后，当段实际负载、实际光伏出力、已经揭示的交易价格以及段初储电量均可获得，储能设备能够按照同一时间粒度完成当段控制；尚未实现的未来负载、光伏和价格不作为实际执行信息。
\item \textbf{信息使用满足因果性}。任一决策时刻的预测、残差场景和参数估计仅使用该时刻以前已经完成的历史观测以及已经正式发布的预报；问题 3 中的 0:00、6:00、12:00、18:00 光伏预报仅在对应发布时间之后进入模型。
\item \textbf{储能效率在研究期内恒定}。按题设将充电效率和放电效率均取 $\eta=0.9$，充电交流侧电量 $C_t$ 对库存增加 $\eta C_t$，放电交流侧电量 $D_t$ 对库存消耗 $D_t/\eta$；题目未提供参数的自放电、电池寿命衰减和启停成本不纳入模型。
\item \textbf{不考虑反向售电收益}。题目仅给出向外网购电和紧急购电的价格机制，未规定微网向外网售电的渠道及价格，因此超过负载和可充电需求的光伏或已购电量允许弃置，且不产生售电收入。
\item \textbf{问题 2--4 按跨日连续库存处理}。除问题 1 按照题目要求施加 $S_0=S_{144}$ 的周期约束外，问题 2--4 不要求每日回到固定库存，当日 24:00 的储电量作为下一日 0:00 的初始状态。为与当前计算结果保持一致，正式评分区间取 2 月 1 日至 12 月 31 日，并在评分期初统一以 6000\,kWh 初始化；1 月数据仅用于预测、场景和参数的因果预热。
\end{enumerate}

紧急购电仅用于补足当段供能缺口，这一处理直接对应题目对紧急购电触发条件的规定，不再作为独立经验假设。同一 10\,min 时段内的充、放电互斥则可由后文互斥消去引理得到，因此无需额外引入二元变量。

\section{符号说明}
记号约定：帽记号 $\hat x$ 表示预测量，上标 $\mathrm{act}$ 表示实际观测，$[z]^+=\max(z,0)$；凸函数 $f$ 的次微分记 $\partial f$。主要符号见表~\ref{tab:symbols}，均在首次出现处另有说明。

\begin{table}[!htbp]\centering\small
\caption{主要符号及适用范围}\label{tab:symbols}
\renewcommand{\arraystretch}{1.08}
\begin{tabularx}{\textwidth}{@{}lXcc@{}}\toprule
\tabheadrow 符号 & 含义 & 单位 & 适用 \\ \midrule
$t,\,T$ & 10 分钟段编号与一天段数，$T=144$ & --- & 一至四 \\
$L_t,\,V_t$ & 段 $t$ 负载、光伏电量 & kWh & 一至四 \\
$p_t$ & 段 $t$ 交易电价（问 1--3 为附件 1 常数曲线） & 元/kWh & 一至四 \\
$G_t,\,G^0_t$ & 购电量；0:00 冻结的计划购电量 & kWh & 一至四 \\
$C_t,\,D_t$ & 储能充、放电量（交流侧） & kWh & 一至四 \\
$S_t$ & 段末储电量，$S_t\in[1200,10800]$ & kWh & 一至四 \\
$\eta,\,\bar P$ & 充放电效率 $0.9$；每段充放电量上限 $5000/6$ & ---, kWh & 一至四 \\
$\psi(x)$ & 库存增量 $x$ 的交流侧电量：$x\ge0$ 为 $x/\eta$，$x<0$ 为 $\eta x$ & kWh & 一至四 \\
$F_t(e)$ & 问 1 确定性剩余最低费用（价值函数） & 元 & 一 \\
$E_t,\,W_t$ & 紧急购电量（5 倍电价）；富余电量 & kWh & 二至四 \\
$M$ & 配对残差场景数，$M=30$ & --- & 二至四 \\
$\omega$ & 场景编号，$\omega=1,\dots,M$ & --- & 二至四 \\
$H_{t,\omega},\,\bar H_t$ & 场景 $\omega$ 未来费用函数；其等权平均 & 元 & 二至四 \\
$R_t$ & 动态保留水平：缺口时不再放电的库存下限 & kWh & 二至四 \\
$V(e)$ & 次日价值函数（次日自由购电的最低费用） & 元 & 二至四 \\
$v$ & 线性续存费率（单日时域配置的末端项系数） & 元/kWh & 二至四 \\
$K_L,\,K_P$ & 负载、光伏历史窗参数（消融配置用） & --- & 二至四 \\
$\hat B_d,\,\hat s_{d,t},\,\beta_d$ & 预测日电量、归一化日内形状、类型切换率 & kWh, ---, --- & 二至四 \\
$\tau,\,q_\tau$ & 发布时刻（0/6/12/18 时）；当天剩余段数 $144-6\tau$ & ---, --- & 三、四 \\
$G^a_t$ & 调整购电量：重优化后对段 $t$ 提交的购电量 & kWh & 三、四 \\
$\Delta^+_t,\,\Delta^-_t$ & 调增、调减量，$G^a=G^0+\Delta^+-\Delta^-$ & kWh & 三、四 \\
$w$ & 0:00 正式预报与历史预测的组合权重 & --- & 三、四 \\
$\widetilde G_{\omega,u}$ & 跨午夜次日暂拟购电（按场景追索，不提交） & kWh & 三、四 \\
$\Gamma_\tau$ & 重优化目标中当天已定购电费与调整费部分 & 元 & 三、四 \\
$\ell_d,\,\chi_{d,t}$ & 日价格水平；归一化日内价格形状 & 元/kWh, --- & 四 \\
$K_p$ & 电价历史窗长度（35 个完整日） & --- & 四 \\
$\lambda$ & 价格同比缩放因子（命题~\ref{prop:scale}） & --- & 四 \\
$K_{\tau,\omega}(e)$ & 场景 $\omega$ 从午夜库存 $e$ 起的跨日续存费用 & 元 & 二至四 \\
\bottomrule\end{tabularx}\end{table}

标签 0:10 表示段 $[0{:}00,0{:}10)$，0:00+1 表示当天最后一段（24:00）。附件 3 的整点预报先插值到段末再转电量。四小时汇总表中充、放电量可同时非零，指该四小时内不同段分别发生充电与放电，同一 10 分钟段内充放电互斥。

% 五、问题一
\section{问题一：确定性计划购电模型}
\subsection{建模思路}
问题一给定一天的电价、负载与光伏，要求在 0:00 一次性确定全天 144 段的购电量与储能充放动作，日初、日末储电量相等，不允许紧急购电。本文先由能量平衡、库存递推与周期约束建立线性规划，获得最优购电费；再证明其价值函数凸、分段线性且边际单调（命题~\ref{prop:convex}），并以凸折线精确表示的连续动态规划独立复算——两条路径的互证为后三问复用同一优化核与折线库奠定基础。

\subsection{确定性线性规划模型}
决策变量为各段购电量 $G_t$、充电量 $C_t$、放电量 $D_t$、富余电量 $W_t$ 与段末储电量 $S_t$，约束来自三方面物理事实，逐条推导如下。

\textbf{能量平衡}：每段供给（购电 + 光伏 + 放电）等于需求（负载 + 充电 + 富余），即
\begin{equation}
G_t+V_t+D_t=L_t+C_t+W_t,\qquad t=1,\dots,144.
\label{eq:q1balance}
\end{equation}
题面“供电不可低于负载”由 $W_t\ge0$ 表达：允许多买，多买部分作为富余弃掉。

\textbf{库存递推与界}：充电入库乘效率 $\eta$、放电出库除以 $\eta$，
\begin{equation}
S_t=S_{t-1}+\eta C_t-D_t/\eta,\qquad 1200\le S_t\le10800,
\label{eq:q1soc}
\end{equation}
且每段充放电量受功率限额约束 $0\le C_t,D_t\le\bar P=5000/6$\,kWh。

\textbf{周期库存}：题面要求 0:00 与 24:00 储电量相同，
\begin{equation}
S_0=S_{144}=6000.
\label{eq:q1periodic}
\end{equation}

目标为全天购电费最小。综合起来得到完整的线性规划模型：
\begin{equation}
\begin{aligned}
\min_{G,C,D,W,S}\quad & \sum_{t=1}^{144}p_tG_t\\
\mathrm{s.t.}\quad
& G_t+V_t+D_t=L_t+C_t+W_t, && t=1,\dots,144,\\
& S_t=S_{t-1}+\eta C_t-D_t/\eta, && t=1,\dots,144,\\
& 1200\le S_t\le 10800,\quad S_0=S_{144}=6000, && \\
& 0\le C_t,\,D_t\le\bar P,\quad G_t,\,W_t\ge 0, && t=1,\dots,144.
\end{aligned}
\label{eq:q1lp}
\end{equation}
目标中不加入 $\varepsilon(C+D+W)$ 一类经济含义之外的正则项；解的退化（同段同时充放电）在动作恢复时按等价变换消除，不影响费用。

\subsection{连续状态动态规划与价值函数性质}
线性规划给出最优值，但只给一条最优轨迹。为了回答“库存偏离该轨迹时未来费用如何变化”——这是后三问执行层的核心——定义剩余最低费用（价值函数）
\begin{equation}
F_t(e)=\min\Bigl\{\text{段 }t\text{ 起的购电费}\;\Big|\;S_{t-1}=e,\ \text{满足式\eqref{eq:q1balance}--\eqref{eq:q1periodic}}\Bigr\},
\label{eq:q1value}
\end{equation}
它满足动态规划\cite{bertsekas}的逆向递推
\begin{equation}
F_t(e)=\min_{x\in\mathcal A_t(e)}\bigl\{p_t\,[L_t-V_t+\psi(x)]^+ + F_{t+1}(e+x)\bigr\},
\qquad
F_{145}(e)=\begin{cases}0,&e=6000,\\ +\infty,&\text{其他},\end{cases}
\label{eq:q1dp}
\end{equation}
其中 $x$ 为库存增量，$\psi(x)$ 为其交流侧电量（表~\ref{tab:symbols}），可行集 $\mathcal A_t(e)$ 由功率限额与库存界给出。

\begin{proposition}[确定性价值函数的凸分段线性结构]\label{prop:convex}
按式\eqref{eq:q1dp} 定义的 $F_t(e)$，$t=1,\dots,145$，是其有效域（$[1200,10800]$ 与终端可达性约束之交）上的凸分段线性函数，且边际价值 $-\partial F_t(e)$ 关于 $e$ 单调非增；断点数至多逐段增加常数个。
\end{proposition}

\begin{proof}
对 $t$ 逆向归纳。$F_{145}$ 为单点指示函数，凸。设 $F_{t+1}$ 凸分段线性。记单段动作费用 $a_t(x)=p_t[L_t-V_t+\psi(x)]^+$。

(i) \emph{$a_t$ 凸分段线性}：$\psi(x)$ 是过原点的两段凸折线（斜率 $\eta<1/\eta$ 递增），$L_t-V_t+\psi(x)$ 仍凸；外层 $p_t[\cdot]^+$ 是非降凸函数与凸函数的复合，故 $a_t$ 凸，断点至多两个（$x=0$ 与 $[\cdot]^+$ 的零点）。

(ii) \emph{一步递推是下确界卷积}：换元 $u=e+x$，
\[
F_t(e)=\min_{u}\{\check a_t(e-u)+F_{t+1}(u)\}=(\check a_t\,\square\,F_{t+1})(e),\qquad \check a_t(z)=a_t(-z),
\]
再与功率限额、库存界对应的区间指示函数相加。凸分段线性函数的下确界卷积仍凸分段线性：其上境图是两上境图的 Minkowski 和，等价于把两条折线的斜率段按斜率归并后首尾相接，断点数为二者之和；与区间指示函数相加即定义域截断，保凸。故 $F_t$ 凸分段线性，断点增量不超过 $a_t$ 的段数（常数）。

(iii) \emph{边际单调}：凸函数的次微分关于自变量单调非减\cite{rockafellar}，故边际价值 $-\partial F_t$ 非增。归纳完成。
\end{proof}


命题~\ref{prop:convex} 的三点含义贯穿全文：(i) $F_t$ 可以用\textbf{断点--斜率序列精确表示}，逆向递推化为凸折线的下确界卷积（斜率归并排序），全程无库存网格、无插值误差；(ii) 边际价值 $-\partial F_t(e)$ 单调非增，“库存越满、追加一度电越不值钱”，这正是储能套利的边际递减结构；(iii) 给定 $F_{t+1}$，最优动作由边际价值与当段电价的比较刻画，即“放电到保留水平为止”的阈值规则——后三问的 DP 价值执行器就是该规则在不完整信息下的直接推广。

图~\ref{fig:fig_f02_value_function} 将该结构可视化：左图为四个代表时刻的 $F_t(e)$ 凸折线，右图为 12:00 的边际价值阶梯。可以看到清晨（$t=36$ 前）边际价值高——此后是早高峰放电窗口；库存越高曲线越平——富余库存的机会成本递减。

\begin{figure}[!htbp]\centering
\safeincludegraphics[width=0.92\textwidth]{fig_f02_value_function}
\caption{问题一价值函数 $F_t(e)$（左）与 12:00 边际价值 $-\partial F_t$（右）：凸、分段线性、边际单调，验证命题~\ref{prop:convex}}
\label{fig:fig_f02_value_function}
\end{figure}

\subsection{双方法互证与求解验证}
先说明二者为何必须相等。式\eqref{eq:q1value} 按定义就是"从状态 $(t,e)$ 出发的原问题最优值"，故 $F_1(6000)$ 与线性规划\eqref{eq:q1lp} 是\emph{同一个}优化问题：可行域相同（能量平衡、库存递推与界、周期约束逐条对应），目标相同。而逆向递推式\eqref{eq:q1dp} 与定义式\eqref{eq:q1value} 的等价即 Bellman 最优性原理\cite{bertsekas}——本问题逐段费用可加、状态转移只依赖当前状态与动作，原理的适用条件全部满足；且凸折线表示对 $F_t$ 是精确的（命题~\ref{prop:convex}），递推中不引入任何离散化或截断近似。因此\textbf{理论上 LP 最优值恒等于 $F_1(6000)$，任何数值差异只可能来自实现错误}——这正是对照的意义：它检验的不是模型，而是两条独立实现路径的代码正确性。

线性规划用 SciPy\cite{scipy} 的 HiGHS 求解器求解；连续动态规划按式\eqref{eq:q1dp} 以凸折线精确递推并正向恢复动作。两条完全独立的实现路径给出同一最优值：
\[
\text{LP}=35126.9486\ \text{元},\qquad F_0(6000)=35126.9486\ \text{元},\qquad \text{差}=7.3\times10^{-12}\ \text{元},
\]
再对非网格初值、零电价、单点终端等 750 组变式做同样对照，最大费用差 $3.64\times10^{-12}$ 元。互证排除了实现错误，也使后三问放心复用同一优化核与折线库。

\subsection{最优策略及经济解读}
最优计划全天购电 59482.70\,kWh、购电费 \textbf{35126.95 元}。作为对照，若不使用储能、逐段按需购电，全天费用为 48052.05 元：储能的峰谷套利把日购电费降低 12925.10 元（\textbf{26.90\%}）。指定时段结果见表~\ref{tab:q1purchase} 与表~\ref{tab:q1storage}，完整策略已存入 result1.xlsx。

\begin{table}[!htbp]
\centering\small
\caption{问题一指定时段购电量（kWh）}\label{tab:q1purchase}
\begin{tabular}{lrrrrrr}
\toprule
时段起点 & 10:00 & 12:00 & 14:00 & 16:00 & 18:00 & 20:00 \\
\midrule
购电量 & 0.00 & 480.41 & 0.00 & 445.43 & 531.89 & 0.00 \\
\bottomrule
\end{tabular}
\end{table}

\begin{table}[!htbp]
\centering\small
\caption{问题一四小时充放电汇总（kWh）}\label{tab:q1storage}
\begin{tabular}{lrr@{\qquad}lrr}
\toprule
时段 & 充电量 & 放电量 & 时段 & 充电量 & 放电量 \\
\midrule
0:00--4:00  & 4500.00 &    0.00 & 4:00--8:00   &  833.33 & 6365.84 \\
8:00--12:00 & 4787.96 & 1703.00 & 12:00--16:00 & 5286.04 &   91.10 \\
16:00--20:00 &   0.00 & 5780.13 & 20:00--24:00 & 5333.33 & 2859.87 \\
\bottomrule
\end{tabular}
\end{table}


\begin{figure}[!htbp]\centering
\safeincludegraphics[width=0.9\textwidth]{fig_f03_q1_day}
\caption{问题一最优调度：(a) 负载、光伏与计划购电功率；(b) 储电量轨迹。低价谷段（0--6 时、12--14 时）满功率充电、购电顶到负载之上，高价峰段（7--11 时、17--20 时）放电、购电归零}
\label{fig:fig_f03_q1_day}
\end{figure}

图~\ref{fig:fig_f03_q1_day} 显示最优策略的经济机理：购电集中在电价谷段并同步充电，峰段完全依靠光伏与放电覆盖负载；储电量在两个峰前充至接近上限、峰内放至接近下限，一天完成约两个完整充放循环——在效率损耗 $\eta^2=0.81$ 下，只有峰谷价差超过约 $1/0.81-1\approx23\%$ 的时段对才值得套利，这解释了为何部分中等价差时段储能保持静默。周期约束的对偶价格给出 24:00 库存的边际影子价，与次日 0:00 的 $F$ 曲线斜率一致，说明周期解自洽。

% 六、问题二
\section{问题二：不确定负载与光伏下的计划购电与因果执行}
问题 2 把负载与光伏变成随机量，模型分四层展开：预测层给出点预测（\S\ref{sec:q2forecast}），场景层把历史残差配对成经验场景（\S\ref{sec:q2scen}），采购层以样本平均近似冻结计划购电量（\S\ref{sec:q2saa}），执行层用 DP 价值执行器逐段因果控制储能（\S\ref{sec:q2dp}）。全部环节严格因果，评价以真实交付数据结算的实际账单为准：
\begin{equation}
C_{\mathrm{real}}=\sum_{d}\sum_{t}\bigl(p_tG^0_{d,t}+5p_tE_{d,t}\bigr).
\label{eq:q2bill}
\end{equation}

\subsection{预测层：分解式负载预测与三日光伏均值}\label{sec:q2forecast}
负载预测的经典做法是类型日均值或统计回归\cite{kangcq,hong}。本题负载呈现明显的双类型结构：周五、周六为低负载日（记 $k=0$），其余为高负载日（$k=1$）；类型由一月前 14 日的平均日电量识别，1 月 15 日起固定。将负载预测分解为\textbf{日电量水平 $\times$ 归一化日内形状}：记 $B_i=\sum_tL_{i,t}$ 为日 $i$ 的日电量，$I_d$ 为 $d$ 之前最近三个同类型日，$J_d$ 为此前 35 天内发生类型切换的日期集，则
\begin{align}
\hat s_{d,t}&=\frac{\sum_{i\in I_d}L_{i,t}}{\sum_{i\in I_d}B_i}, &&\text{（同类型日内形状）}\label{eq:q2shape}\\
\beta_d&=\operatorname*{median}_{i\in J_d}\frac{\ln(B_i/B_{i-1})}{k_i-k_{i-1}}, &&\text{（类型切换的对数电量比率）}\label{eq:q2beta}\\
\hat B_d&=B_{d-1}\exp\bigl[\beta_d(k_d-k_{d-1})\bigr],\qquad
\hat L_{d,t}=\hat B_d\,\hat s_{d,t}. &&\text{（水平递推 $\times$ 形状）}\label{eq:q2load}
\end{align}
预测次日时，水平递推使用当日的\emph{预测}日电量而非尚未结束的真实总量。1 月 15 日前采用当时可得的同星期均值，样本不足回退至最近至多七日均值。光伏点预测取此前最多三个完整日的逐段均值（光伏功率预测的系统综述见文献\cite{antonanzas}；本题的正式数值预报在问题 3 才可用）。图~\ref{fig:fig_f04_load_forecast} 给出 3 月 20 日的预测效果：形状复现准确，残差集中在光伏午间波动段——这正是场景层要刻画的不确定性。

\begin{figure}[!htbp]\centering
\safeincludegraphics[width=0.9\textwidth]{fig_f04_load_forecast}
\caption{3 月 20 日 0:00 负载预测与实际（上）及残差（下）：分解式预测复现日内双峰形状，残差无系统偏倚}
\label{fig:fig_f04_load_forecast}
\end{figure}

\subsection{场景层：完整路径配对残差}\label{sec:q2scen}
预测误差不是逐段独立噪声——负载误差日内强相关，光伏误差与天气过程绑定，两通道之间也相关。因此场景不做参数化假设，而直接使用\textbf{完整路径残差}：每条历史 0:00 发布预测在其 24 小时全部交付后，与实际曲线相减得到一条负载残差与一条光伏残差，二者来自同一历史日、保持配对。决策日 0:00 选取此前最近 $M=30$ 条完整配对残差，等权叠加到当天点预测（负值截零），得到场景集
\begin{equation}
L_{\omega,t}=\bigl[\hat L_{d,t}+r^{L}_{\omega,t}\bigr]^+,\qquad
V_{\omega,t}=\bigl[\hat V_{d,t}+r^{V}_{\omega,t}\bigr]^+,\qquad \omega=1,\dots,M.
\label{eq:q2scenario}
\end{equation}
历史残差必须由“历史当时可得信息”生成的预测计算（一月起因果重建），未完成的路径不入库。图~\ref{fig:fig_f05_scenario_fan} 显示 30 条净负荷场景对真值的覆盖：场景带在光伏时段最宽，与不确定性来源一致。

\begin{figure}[!htbp]\centering
\safeincludegraphics[width=0.9\textwidth]{fig_f05_scenario_fan}
\caption{3 月 20 日净负荷（负载减光伏）的 30 条配对场景与真值：场景带覆盖真值路径，午间光伏波动段最宽}
\label{fig:fig_f05_scenario_fan}
\end{figure}

\subsection{采购层：样本平均近似与 48 小时主方法}\label{sec:q2saa}
0:00 决策“买多少”建模为两阶段随机规划\cite{saa-book,birge}的样本平均近似（SAA）\cite{kleywegt}：第一阶段变量为全天计划购电量 $G^0$（跨场景共享），第二阶段为各场景的储能追索与缺口补购 $(C_\omega,D_\omega,S_\omega,E_\omega,W_\omega)$。约束逐条与式\eqref{eq:q1lp} 同构，仅平衡式换成场景曲线并加入紧急购电项。完整模型为
\begin{equation}
\begin{aligned}
\min_{G^0,\{C_\omega,D_\omega,S_\omega,E_\omega,W_\omega\}}\quad
& \sum_{t}p_tG^0_t+\frac1M\sum_{\omega=1}^{M}\Bigl(\sum_{t}5p_tE_{\omega,t}-vS_{\omega,144}\Bigr)\\
\mathrm{s.t.}\quad
& G^0_t+E_{\omega,t}+V_{\omega,t}+D_{\omega,t}=L_{\omega,t}+C_{\omega,t}+W_{\omega,t}, && \forall\omega,t,\\
& S_{\omega,t}=S_{\omega,t-1}+\eta C_{\omega,t}-D_{\omega,t}/\eta,\quad S_{\omega,0}=S_{\mathrm{now}}, && \forall\omega,t,\\
& 1200\le S_{\omega,t}\le10800,\quad 0\le C_{\omega,t},D_{\omega,t}\le\bar P, && \forall\omega,t,\\
& G^0_t,\,E_{\omega,t},\,W_{\omega,t}\ge0, && \forall\omega,t.
\end{aligned}
\label{eq:q2saa}
\end{equation}
末端项 $-vS_{\omega,144}$ 给 24:00 结转库存定价。24 小时基线配置取线性续存费率 $v=0.481548$ 元/kWh——即次日 0--5 时均价 $0.4334$ 除以放电效率 $\eta$：结转一度库存可替代次日凌晨 $1/\eta$ 度购电。\textbf{48h+$V(e)$ 主方法}把时域延伸到次日：次日按共用点预测曲线自由购电、末端价值为零，等价于用次日价值函数 $V(e)$（次日自由购电的确定性最低费用，按问 1 的凸折线递推精确构造）替代线性项 $-ve$，从而以跨日机会成本为午夜库存定价。

需要明确该目标的性质：第二阶段按场景独立追索、场景内完美预见，故式\eqref{eq:q2saa} 是多阶段问题的乐观代理（wait-and-see 近似）。它只用来定 $G^0$；真实执行不用任何场景内动作，而由下面的执行层因果生成。

\subsection{执行层：未来费用函数与 DP 价值执行器}\label{sec:q2dp}
$G^0$ 冻结后，执行层唯一的自由度是每段的充放电。困难在于放电有机会成本：现在放掉的库存，可能在未来更贵的缺口段更值钱。为此把问 1 的价值函数推广到固定采购、场景平均的形式。对每条场景 $\omega$ 逆向递推
\begin{equation}
H_{u,\omega}(e)=\min_{x}\bigl\{\phi_{u,\omega}(x)+H_{u+1,\omega}(e+x)\bigr\},\qquad
\bar H_u(e)=\frac1M\sum_{\omega=1}^{M}H_{u,\omega}(e),
\label{eq:q2hbar}
\end{equation}
其中当天固定采购段记净缺口 $r=L_{\omega,u}-V_{\omega,u}-G^0_u$，单段费用 $\phi_{u,\omega}(x)=5p_u[r+\psi(x)]^+$（缺口只能放电或 5 倍电价补购，富余只允许充电）；跨午夜自由采购段 $\phi_{u,\omega}(x)=p_u[L-V+\psi(x)]^+$；末端接 $V(e)$（主方法）或 $-ve$（基线）。由命题~\ref{prop:convex} 的同一论证，每个 $H_{u,\omega}$ 与 $\bar H_u$ 均为凸折线，可精确表示与平均。

\textbf{因果执行规则（DP 价值执行器）}按下述步骤逐段进行。

\textbf{Step 1}（观测）：段 $t$ 开始时观测真实负载 $L_t^{\mathrm{act}}$、光伏 $V_t^{\mathrm{act}}$ 与电价，计算净缺口 $r_t=L_t^{\mathrm{act}}-V_t^{\mathrm{act}}-G^0_t$。

\textbf{Step 2}（富余充电）：若 $r_t\le0$，充电 $C_t=\min\{-r_t,\ \bar P,\ (10800-S_{t-1})/\eta\}$，$D_t=E_t=0$，转 Step 5。

\textbf{Step 3}（保留水平）：若 $r_t>0$，由次段平均未来费用与当段紧急电价计算动态保留水平
\begin{equation}
R_t=\max\operatorname*{arg\,min}_{1200\le z\le10800}\bigl\{5p_t\eta z+\bar H_{t+1}(z)\bigr\}.
\label{eq:q2reserve}
\end{equation}
其含义：库存中边际价值超过当段紧急购电边际代价 $5p_t\eta$ 的部分应当留存，$R_t$ 正是该边际比较的平衡点；由 $\bar H_{t+1}$ 凸性，式\eqref{eq:q2reserve} 在折线断点处即可精确求得。

\textbf{Step 4}（限额放电）：放电 $D_t=\min\{r_t,\ \bar P,\ \eta[S_{t-1}-R_t]^+\}$，剩余缺口紧急购电 $E_t=r_t-D_t$。

\textbf{Step 5}（结转）：$S_t=S_{t-1}+\eta C_t-D_t/\eta$，富余 $W_t$ 由平衡式回填；进入下一段。

该规则只使用当段已观测量与 0:00 已构造的 $\bar H$，天然互斥充放电，且紧急购电从不用于充电（假设 5）。图~\ref{fig:fig_f06_dp_mechanism} 展示执行机理：保留水平随时段起伏——未来有高价缺口时抬高（惜放），临近日末且次日凌晨便宜时压低（敢放）；SOC 轨迹在 $R_t$ 之上时缺口全由放电覆盖，触及 $R_t$ 后转为紧急购电。

\begin{figure}[!htbp]\centering
\safeincludegraphics[width=0.92\textwidth]{fig_f06_dp_mechanism}
\caption{DP 价值执行器机理（3 月 20 日）：上图为动态保留水平 $R_t$ 与储电量轨迹及紧急购电段；下图为三个时刻的 $\bar H_t(e)$ 折线——保留水平即“边际价值 = 当段紧急代价”的平衡点}
\label{fig:fig_f06_dp_mechanism}
\end{figure}

\subsection{候选计划重评}
SAA 解 $G^s$ 优化的是完美追索代理，与因果执行之间存在口径差。为弥合该差距，引入候选计划重评：取点预测解 $G^d$（零残差单场景 LP 的购电量），构造凸组合候选 $G(a)=(1-a)G^s+aG^d$，$a\in\{0,0.25,0.5,0.75,1\}$；每个候选独立重建 $\bar H$、沿同批场景按上述因果规则模拟执行，以模拟账单（含各场景自身的跨午夜续存费用 $K_{\tau,\omega}$，保持路径配对）评分，取最低者冻结为 $G^0$。评分时把 $K_{\tau,\omega}$ 换成平均 $\bar K$ 会破坏跨午夜相关性，属于实现禁忌。

\subsection{结果与解读}
主方法（48h+$V(e)$）全年（2 月 1 日--12 月 31 日，334 个评分日）实际费用 \textbf{1352.96 万元}，其中紧急购电费 78.23 万元（购电量占比 0.69\%）；24 小时基线为 1356.64 万元，时域与终端定价组件全年净省 3.68 万元。定位区间：完美信息离线下界（全年真值一次 LP、无紧急与调整费）为 1222.82 万元，无储能直购（已知真值按需购电）为 1640.73 万元，与主方法同一预测下无储能定购为 1879.64 万元。主方法费用位于下界之上 10.6\%、比同信息无储能方案节省 28.0\%——差距的主要来源是信息不完整（预测残差引发的紧急购电与保守购电），而非执行规则本身。

\begin{table}[!htbp]
\centering\small
\caption{24小时基准配置：334个评分日的实际账单与末库存}\label{tab:baseline}
\begin{tabular}{|l|r|r|r|r|}
\hline
\tabheadrow 分支 & 总费/万元 & 紧急费/万元 & 紧急量/kWh & 末库存/kWh \\ \hline
问2 & 1356.6396 & 78.3421 & 143647.07 & 9459.10 \\ \hline
问3 & 1302.4597 & 19.2075 & 39811.75 & 9623.58 \\ \hline
问4-2 & 1427.0607 & 78.6925 & 137553.82 & 9455.36 \\ \hline
问4-3 & 1371.6972 & 19.3585 & 38118.94 & 9635.83 \\ \hline
\end{tabular}
\end{table}

\begin{table}[!htbp]
\centering\small
\caption{跨日前瞻配置（48h+$V(e)$）：334个评分日的实际账单与末库存}\label{tab:main-results}
\begin{tabular}{|l|r|r|r|r|}
\hline
\tabheadrow 分支 & 总费/万元 & 紧急费/万元 & 紧急量/kWh & 末库存/kWh \\ \hline
问2 & 1352.9635 & 78.2275 & 143233.08 & 8349.65 \\ \hline
问3 & 1298.9551 & 19.1602 & 39569.80 & 8571.66 \\ \hline
问4-2 & 1422.9900 & 79.2771 & 139099.84 & 9329.24 \\ \hline
问4-3 & 1369.3012 & 19.5048 & 38269.94 & 9024.54 \\ \hline
\end{tabular}
\end{table}


\begin{figure}[!htbp]\centering
\safeincludegraphics[width=0.92\textwidth]{fig_f07_q2_monthly}
\caption{问题二主方法月度费用与参照带：12 月最高（冬季负载高、光伏弱），6--7 月次之（空调负载），4--5 月春季最低；全年每月都稳定落在“下界—无储能”区间内侧}
\label{fig:fig_f07_q2_monthly}
\end{figure}

月度分解（图~\ref{fig:fig_f07_q2_monthly}）显示费用的季节结构：12 月最贵（165.68 万元，冬季负载高、光伏出力弱），6--7 月为次高峰（137.65 与 144.79 万元，空调负载抬升日电量），4--5 月最便宜（94.99--97.76 万元，光伏强而负载温和）。指定日期结果见 \S\ref{sec:reportdays} 统一给出。

% 七、问题三
\section{问题三：正式预报与滚动调整购电}
问题 3 在问题 2 之上扩展两件事：\textbf{信息}——每天 0:00、6:00、12:00、18:00 获得未来 24 小时整点光伏预报；\textbf{权限}——可在发布时刻调整尚未执行的购电量，调减部分付 50\% 违约费、调增部分按 1.5 倍电价购入。全天实际账单为四项分解
\begin{equation}
C_{\mathrm{real}}=\sum_t\Bigl[p_t\min(G^0_t,G^a_t)+0.5\,p_t\,(G^0_t-G^a_t)^+
+1.5\,p_t\,(G^a_t-G^0_t)^++5p_tE_t\Bigr],
\label{eq:q3bill}
\end{equation}
等价于 $\sum_tp_t[G^a_t+0.5|G^a_t-G^0_t|+5E_t]$：交付按最终执行量计价，偏离原计划的部分对称地付半价摩擦费。每个交付段只按最终提交量结算一次，被覆盖的中间临时决策不计费。

\subsection{预报接入与组合预测}
附件 3 发布后第 $h$ 小时的预报对应发布时刻加 $h$ 的整点值。以发布时刻已观测光伏为左锚点，把 24 小时整点预报插值到 10 分钟段末。0:00 的正式预报与历史预测（三日均值）各有所长，取凸组合
\begin{equation}
\widehat{V}=w\,\widehat{V}^{F}+(1-w)\,\widehat{V}^{H},\qquad w=0.388815,
\label{eq:q3combine}
\end{equation}
其中权重 $w$ 为最小化组合误差平方和的经验权重 $w=\sum e_H(e_H-e_F)/\sum(e_F-e_H)^2$（截断到 $[0,1]$），由一月内各历史发布时刻此前已完成的配对误差因果估计，2 月 1 日起冻结。6/12/18 时的预报误差显著更小，直接采用最新正式预报。滚动场景库按发布小时分组：每次只选\emph{同发布小时}、完整 24 小时已实现的最近 30 条三通道配对残差，跨午夜路径按精确结束时刻判断可用性。

\subsection{滚动重优化模型}
发布时刻 $\tau$（当天剩余 $q_\tau=144-6\tau$ 段）的重优化以当前真实储电量为初值，在冻结的 $G^0$ 上决策调整量。第一阶段变量为当天剩余段的调增 $\Delta^+_u$ 与调减 $\Delta^-_u$（跨场景共享），$G^a_u=G^0_u+\Delta^+_u-\Delta^-_u$；第二阶段为各场景储能追索、缺口补购与跨午夜次日暂拟购电 $\widetilde G_{\omega,u}$（按场景分开，只用于估计未来费用，不提交）。完整模型：
\begin{equation}
\begin{aligned}
\min_{\Delta^\pm,\{\cdot_\omega\}}\quad
& \Gamma_\tau+\sum_{u>\!6\tau}\bigl(1.5p_u\Delta^+_u-0.5p_u\Delta^-_u\bigr)
+\frac1M\sum_{\omega}\Bigl[\sum_{u\in\text{次日}}p_u\widetilde G_{\omega,u}
+\sum_{u}5p_uE_{\omega,u}\Bigr]-\frac1M\sum_\omega V(S_{\omega,\mathrm{end}})\\
\mathrm{s.t.}\quad
& G^a_u=G^0_u+\Delta^+_u-\Delta^-_u,\quad 0\le\Delta^-_u\le G^0_u,\ \Delta^+_u\ge0, && u>6\tau,\\
& G^a_u+E_{\omega,u}+V_{\omega,u}+D_{\omega,u}=L_{\omega,u}+C_{\omega,u}+W_{\omega,u}, && \forall\omega,\ u>6\tau,\\
& \widetilde G_{\omega,u}+E_{\omega,u}+V_{\omega,u}+D_{\omega,u}=L_{\omega,u}+C_{\omega,u}+W_{\omega,u}, && \forall\omega,\ u\in\text{次日},\\
& S_{\omega,u}=S_{\omega,u-1}+\eta C_{\omega,u}-D_{\omega,u}/\eta,\quad S_{\omega,6\tau}=S_{\mathrm{now}}, && \forall\omega,u,\\
& 1200\le S_{\omega,u}\le10800,\ 0\le C_{\omega,u},D_{\omega,u}\le\bar P,\ E,W,\widetilde G\ge0, && \forall\omega,u,
\end{aligned}
\label{eq:q3lp}
\end{equation}
其中 $\Gamma_\tau$ 为按式\eqref{eq:q3bill} 结算的当天已执行部分费用（常数），$V(\cdot)$ 为次日价值函数（主方法；24h 基线用 $-vS$）。调增与调减在同段的净费用为 $+p_u>0$，故最优解中二者不共存。前瞻始终为完整 24 小时（如 18:00 覆盖当天 18--24 时与次日 0--18 时），但\textbf{只提交到下一发布时刻}：之后的段是临时决策，下轮重优化以更新的信息覆盖之。执行层与问题 2 完全相同——每轮重优化后按新的 $\bar H$ 与保留水平因果执行到下一发布时刻。

\begin{proposition}[重优化的代理不劣性]\label{prop:readjust}
固定发布时刻 $\tau$ 的信息状态（同一场景集、同一未来预测、同一当前储电量与冻结计划 $G^0$）。式\eqref{eq:q3lp} 的重优化线性规划的最优值不超过不调整方案（$\Delta^\pm\equiv0$，即 $G^a=G^0$）的代理费用；差值即该发布时刻在代理口径下的信息价值，非负。
\end{proposition}

\begin{proof}
取 $\Delta^+=\Delta^-=0$，其余变量取不调整方案的第二阶段最优追索值，则该组取值满足式\eqref{eq:q3lp} 的全部约束，是其可行解，目标值恰为不调整方案的代理费用。线性规划最优值不超过任一可行解的目标值，命题得证。
\end{proof}


命题~\ref{prop:readjust} 给出的是\emph{代理口径}的方向保证：更多发布时刻不会让重优化 LP 变差。实际账单由真值结算，不被该命题覆盖，需由全年实验回答（\S\ref{sec:q3issues}）。数值上，四个指定日 12:00 重优化的代理改善最小值为 56.70 元 $\ge0$，与命题一致。

图~\ref{fig:fig_f08_rolling_timeline} 给出滚动机制的时间结构：每轮 24 小时前瞻、只提交一个发布间隔、临时决策可被覆盖。

\begin{figure}[!htbp]\centering
\safeincludegraphics[width=0.94\textwidth]{fig_f08_rolling_timeline}
\caption{滚动重优化时间线：四个发布时刻各前瞻完整 24 小时（可跨午夜），实际只提交到下一发布时刻（实色段），其后为可被覆盖的临时决策（虚框段）}
\label{fig:fig_f08_rolling_timeline}
\end{figure}

\subsection{是否需要引入其他时刻的预报}\label{sec:q3issues}
固定预测方法、场景构造、终端定价与结算口径，逐级开放发布时刻集合 $\{0\}\subset\{0,6\}\subset\{0,6,12\}\subset\{0,6,12,18\}$，四组配置均连续执行全年：

\begin{table}[!htbp]
\centering\small
\caption{问3发布时刻组合与相邻组合的费用变化}\label{tab:issues}
\begin{tabular}{lrrr}
\toprule
发布时刻/h & 总费用/元 & 紧急费/元 & 相邻节省/元 \\
\midrule
0 & 13351182.68 & 579245.64 & --- \\
0,6 & 13312853.41 & 353878.30 & 38329.27 \\
0,6,12 & 13075459.34 & 265785.37 & 237394.07 \\
0,6,12,18 & 13024596.85 & 192075.49 & 50862.49 \\
\bottomrule
\end{tabular}
\end{table}


相邻差即该发布时刻的边际价值（在此固定增加顺序下）：新增 6:00 全年节省 3.83 万元，新增 \textbf{12:00 节省 23.74 万元}，新增 18:00 节省 5.10 万元（图~\ref{fig:fig_f09_issue_marginal}）。12:00 的边际贡献为何一枝独秀？两个因素在此交汇：其一，光伏不确定性集中在午间出力峰值段，12:00 预报恰在峰值前夕把当天剩余光伏路径的方差大幅压缩；其二，12:00 时当天还有一半的段未执行，调整窗口足够长，信息可以兑换成实际的采购修正。相比之下，6:00 预报虽早，但此时距光伏峰值尚远、预报本身的增量信息少；18:00 预报虽准，但当天只剩 6 小时可调，信息变现空间小。\textbf{结论：需要引入其他时刻的预报，其中 12:00 预报的引入最为关键}；若受成本限制只能保留一个日内更新时刻，应保留 12:00。

\begin{figure}[!htbp]\centering
\safeincludegraphics[width=0.86\textwidth]{fig_f09_issue_marginal}
\caption{逐级增加发布时刻的全年边际节省：12:00 预报的边际价值（23.74 万元/年）约为 6:00 与 18:00 之和的三倍——信息增量与剩余调整窗口在午间交汇}
\label{fig:fig_f09_issue_marginal}
\end{figure}

\subsection{结果}
主方法（48h+$V(e)$）全年实际费用 \textbf{1298.96 万元}，比问题 2 主方法降低 54.01 万元（4.0\%）：这是正式预报信息与日内调整权限的联合价值。四项费用分解中紧急购电费降至 19.16 万元（问题 2 为 78.23 万元）——日内调整把大部分本须紧急补购的缺口提前转化为 1.5 倍价的调增，避开了 5 倍价；紧急购电量占比从 0.70\% 降至 0.17\%。24 小时基线为 1302.46 万元，与问题 2 一致地支持 48h 组件。月度与指定日结果见 \S\ref{sec:results}。

% 八、问题四
\section{问题四：波动电价下的购电策略}
问题 4 把电价变成随机过程，储能的角色从平滑负载扩展为价格套利\cite{hejazi}。附件 4 电价的全年逐段均值恰为附件 1 的常数电价曲线，即问题 2/3 是问题 4 的“平均电价”特例；波动电价一方面放大峰谷价差带来的套利空间，另一方面让“何时买、何时调”的判断依赖于\emph{还未实现}的价格。模型在前两问框架上只新增一个电价预测层，并按交易权限分为两支重算。

\subsection{电价预测层：日水平 $\times$ 日内形状}
电价预测文献\cite{weron}表明日内形状远比日间水平易于预测。电价与净负荷（负载减光伏）同受供需驱动，采用与负载预测同构的分解结构。对此前 $K_p=35$ 个完整日计算日价格水平 $\ell_i=T^{-1}\sum_tp_{i,t}$ 与平均归一化日内形状 $\chi_{d,t}$，两支分别建立水平模型：

\textbf{4-2 支（净负荷回归）}：价格水平由前一日水平与当日净负荷预测回归，
\begin{equation}
\hat\ell_{d\mid\tau}=a+\rho\,\ell_{d-1}+b\,\widehat N_{d\mid\tau}/10^5,\qquad
\hat p_{d,t\mid\tau}=\hat\ell_{d\mid\tau}\,\chi_{d,t},
\label{eq:q4price}
\end{equation}
其中 $\widehat N_{d\mid\tau}$ 为发布时刻 $\tau$ 的当日净负荷估计（已观测段真值 + 剩余段预测）。关键的因果细节：回归的\emph{历史}输入同样使用“历史当时的净负荷预测”重建，绝不用历史真实全天净负荷替代——否则拟合环境与使用环境不一致。\textbf{4-3 支}使用不含净负荷项的 AR(1) 水平模型。日内更新时按最近已完成段的预测误差做衰减修正；样本不足或回归退化按因果回退规则处理，预测价格下限 $10^{-4}$。

价格残差与负载、光伏残差按\textbf{同一历史发布时刻的完整路径}三通道配对（式\eqref{eq:q2scenario} 增加价格通道），保持“高负载日往往高电价”的跨通道相关。优化目标中：调整与普通购电费用（跨场景共享的第一阶段变量）按场景均价计费，各场景的紧急购电与跨午夜暂拟购电按该场景自身价格计费，实际结算始终使用附件 4 真值。

图~\ref{fig:fig_f10_price_forecast} 给出预测质量的两端：3 月 20 日（MAE 0.030 元/kWh）形状与水平俱佳；12 月 21 日（MAE 0.191 元/kWh）遭遇价格水平跳升，回归无法预见。全年逐日 MAE 的这种右偏分布正是问题 4 相对完美价格信息的费用差来源。

\begin{figure}[!htbp]\centering
\safeincludegraphics[width=0.92\textwidth]{fig_f10_price_forecast}
\caption{电价 0:00 预测与真值：左为 3 月 20 日（MAE 0.030 元/kWh，形状水平俱准），右为 12 月 21 日（MAE 0.191 元/kWh，水平跳升不可预见）——诚实展示预测的两端}
\label{fig:fig_f10_price_forecast}
\end{figure}

\subsection{两支的优化模型与权限差异}
\textbf{4-2（对应问题 2 权限）}：仅 0:00 冻结当天采购，不读取附件 3。其 0:00 优化模型与式\eqref{eq:q2saa} 同构，仅把常数电价换成价格场景：
\begin{equation}
\min_{G^0,\{\cdot_\omega\}}\ \sum_t\bar p_tG^0_t
+\frac1M\sum_\omega\Bigl(\sum_t5p_{\omega,t}E_{\omega,t}-v_\tau S_{\omega,144}\Bigr)
\quad\text{s.t. 式\eqref{eq:q2saa} 的全部约束},
\label{eq:q42}
\end{equation}
$\bar p_t$ 为场景均价，$v_\tau$ 为价格场景次日 0--5 时均价除以 $\eta$（随场景刷新，取代常数 0.481548）。虽无交易权限，但\textbf{价值更新不需要权限}：0/6/12/18 时仍重建价格预测、价格场景与 $\bar H$，执行层的保留水平式\eqref{eq:q2reserve} 使用当段\emph{真值}电价现算——价格信息随时间改善，执行随之变聪明，采购保持冻结。

\textbf{4-3（对应问题 3 权限）}：完整沿用式\eqref{eq:q3lp}，把各处电价换成对应的价格场景（调整费用按场景均价、场景内费用按各自价格），每个发布时刻同时刷新光伏预报、价格预测与剩余采购。

\begin{proposition}[价格同比缩放不变性]\label{prop:scale}
设全部价格输入同乘 $\lambda>0$：历史与真实电价 $p\mapsto\lambda p$（紧急电价保持 5 倍关系、调整费保持 $1.5$ 与 $0.5$ 倍关系），续存费率 $v\mapsto\lambda v$。则 (a) 各采购线性规划（问 1 确定性 LP、问 2 SAA、问 3 重优化）的最优解集不变、最优费用同乘 $\lambda$；(b) 价格预测层输出同乘 $\lambda$；(c) 未来费用函数 $\bar H_t\mapsto\lambda\bar H_t$，动态保留水平 $R_t$ 不变，DP 价值执行器的因果动作逐段不变；(d) 实际账单同乘 $\lambda$。即策略形状只由价格的相对结构决定，与价格水平无关。
\end{proposition}

\begin{proof}
(a) 三个线性规划的约束（能量平衡、库存递推、界）均不含价格；价格只出现在目标向量中，且各项费用（普通购电 $p$、调增 $1.5p$、调减 $-0.5p$、紧急 $5p$、续存 $-v$）对 $\lambda$ 一次齐次，故目标向量整体乘 $\lambda$，最优解集不变、最优值乘 $\lambda$。

(b) 价格预测为近 $K_p$ 个完整历史日的水平均值与归一化形状之积再乘水平因子（式\eqref{eq:q4price}）；水平、形状归一化与回归修正各环节均为历史价格的一次齐次函数，故输出乘 $\lambda$。

(c) $\bar H_t$ 由逐段下确界卷积生成：单段动作费用（固定采购段 $5p[r+\psi(x)]^+$、自由段 $p[L-V+\psi(x)]^+$）与终端项 $-\lambda ve$ 均对 $\lambda$ 一次齐次；下确界卷积与逐点平均都与正常数因子可交换，逆向归纳得 $\bar H_t\mapsto\lambda\bar H_t$。于是保留水平目标 $\lambda[\bar H_{t+1}(z)+5p_t\eta z]$ 的最小点集不变，$R_t$ 不变，执行动作逐段不变。

(d) 动作不变、结算价格乘 $\lambda$，账单各项一次齐次，总费用乘 $\lambda$。
\end{proof}


命题~\ref{prop:scale} 的数值验证（全局 0.5$\times$/2$\times$ 缩放）四项全部通过：LP 费用同比且原最优解仍属缩放后最优解集；SAA 计划购电量逐段不变；保留水平逐段不变；电价预测逐段同比。其含义是：问题 4 策略的成败取决于价格的\emph{形状与可预测性}，而非价格水平本身——把全年电价整体抬高一倍不改变任何决策，改变决策的是“今天下午比明天凌晨贵多少、这个差能否被预见”。

\subsection{结果与解读}
主方法全年实际费用：\textbf{4-2 为 1422.99 万元，4-3 为 1369.30 万元}（24h 基线分别为 1427.06 与 1371.70 万元）。三个定位参照：附件 4 电价下的完美信息离线下界为 1278.09 万元；同一算法、完美价格信息（0:00 已知当天与次日真值电价、无电价残差场景，与 24h 基线同口径）的参照为 4-2 支 1418.18 万元、4-3 支 1362.34 万元。两个观察值得展开：

\textbf{其一，价格预测误差的代价意外地小}。同口径比较，4-2 基线与其完美价格参照差 8.88 万元（0.62\%），4-3 基线差 9.36 万元（0.68\%）。原因在结构上：采购与执行依赖的是价格的\emph{相对形状}（何时便宜何时贵），而日内形状 $\chi_{d,t}$ 高度稳定、易于预测；难以预测的是日间水平跳动，而命题~\ref{prop:scale} 恰说明水平误差对决策形状几乎无影响，只影响对少数极端日的调整幅度。

\textbf{其二，4-3 相对 4-2 的 53.69 万元差距与问题 3 相对问题 2 的 54.01 万元几乎一致}。两个差距都主要来自“正式预报 + 日内调整权限”这一组合，说明该组合的价值对电价环境稳健；波动电价没有放大也没有侵蚀调整权限的价值。逐月看（图~\ref{fig:fig_f11_q4_monthly}），4-3 的优势集中在光伏波动大的月份，2 月甚至出现小幅负值——当月负载水平低、调整机会少，而 4-3 的 AR(1) 价格模型在当月水平跳变日恰好方向出错；诚实呈现该负值月份正说明差距的来源是信息而非口径。

\begin{figure}[!htbp]\centering
\safeincludegraphics[width=0.92\textwidth]{fig_f11_q4_monthly}
\caption{问题四两支相对各自常数电价对应问（4-2 对问 2、4-3 对问 3）的逐月费用差：波动电价整体抬升费用约 5\%，抬升幅度冬季最大；4-3 对 4-2 的月度优势与问 3 对问 2 的模式一致}
\label{fig:fig_f11_q4_monthly}
\end{figure}

四个分支处于不同的信息与交易环境，总费用只在各自环境内评价；4-2 与 4-3 之间同时改变了预测规则、可得光伏信息与采购权限，其差额是组合价值，本文不对其做单因素归因。完整策略已存入 result4-2.xlsx 与 result4-3.xlsx。

% 十、模型评价与推广
\section{模型评价与推广}
\subsection{模型的优点}
\begin{enumerate}
\item \textbf{四问一体的统一框架}：同一优化核、同一价值函数链（式\eqref{eq:chain}）、同一执行器贯穿四问，问题间只改变信息状态与交易权限；每个新增组件（正式预报、调整权限、价格场景）都有对应的消融或对照量化其价值，模型的递进是可验证的递进。
\item \textbf{理论结构支撑}：命题~\ref{prop:convex} 使价值函数以凸折线精确表示（无库存网格、无插值误差），命题~\ref{prop:readjust} 保证重优化方向正确，命题~\ref{prop:scale} 说明策略对价格水平结构稳健；三个命题均配独立数值验证。
\item \textbf{严格因果与无信息泄漏}：预测、场景、调参、执行各环节只用当时可得数据；全部参数由一月选定，评分期从未参与调参。
\item \textbf{结果可信度多重锚定}：LP 与连续 DP 双方法互证（750 组对照最大差 $3.64\times10^{-12}$ 元）；143 项自动化测试；全部 22 组全年实验由轨迹独立重算账单，库存递推、能量平衡与费用分解的复核误差均在 $10^{-12}$ 元量级；实验保存输入数据与代码指纹，可精确复现。
\item \textbf{按实际账单评价}：所有结论以真实交付数据结算的连续执行账单为准，并同时报告紧急购电与期末库存，代理目标与真实费用从不混淆。
\end{enumerate}

\subsection{模型的缺点与局限}
\begin{enumerate}
\item SAA 的场景内完美追索是乐观代理，本文未给出多阶段随机最优性的理论保证；平均未来费用函数 $\bar H$ 同样是机会成本近似而非精确随机价值函数。
\item 全部结论基于单一评分年（2025 年）数据，缺少跨年份外部验证；单因子消融收益不保证跨环境可加。
\item 经验残差场景假设“未来误差分布与近期历史相似”，对突变天气或负载结构变化的适应有滞后；场景等权而未做重要性加权。
\item 电价预测的日间水平模型较简单（回归 / AR(1)），极端价格跳变日（如 12 月 21 日）无法预见；所幸命题~\ref{prop:scale} 表明水平误差对决策形状影响有限。
\end{enumerate}

\subsection{改进方向与推广}
改进方向：(i) 对场景做基于天气类型的条件化选择或重要性加权，替代纯时间近邻；(ii) 价格水平模型引入外生温度/节假日特征；(iii) 用逼近动态规划在 $\bar H$ 上叠加误差修正项，压缩乐观代理与真实账单之间的口径差。推广方面，“信息状态—未来费用函数—因果执行”的框架不依赖本题细节：任何“提前承诺 + 库存缓冲 + 事后结算”的运营问题——电动车充电站的日前购电、冷库/蓄热系统的错峰用能、云计算的预留实例采购——都可以在同一骨架下实例化，只需替换预测层与结算式。凸折线价值函数的精确表示对一维库存状态普遍适用，计算量与时段数线性增长，可直接用于更细粒度（如 1 分钟）的执行控制。


\section*{AI 工具使用声明}
本参赛队在竞赛过程中使用了 AI 工具，主要用于代码开发辅助、数据分析与可视化辅助、论文语言润色，详细使用情况见支撑材料。

\begin{thebibliography}{15}
\bibitem{problem} 2026年高教社杯全国大学生数学建模竞赛C题及附件1--4[Z]. 2026.
\bibitem{wangcs} 王成山, 武震, 李鹏. 微电网关键技术研究[J]. 电工技术学报, 2014, 29(2): 1--12.
\bibitem{liuyx} 刘一欣, 郭力, 王成山. 微电网两阶段鲁棒优化经济调度方法[J]. 中国电机工程学报, 2018, 38(14): 4013--4022.
\bibitem{kangcq} 康重庆, 夏清, 张伯明. 电力系统负荷预测研究综述与发展方向的探讨[J]. 电力系统自动化, 2004, 28(17): 1--11.
\bibitem{saa-book} Shapiro A, Dentcheva D, Ruszczy\'nski A. Lectures on Stochastic Programming: Modeling and Theory[M]. Philadelphia: SIAM, 2009.
\bibitem{birge} Birge J R, Louveaux F. Introduction to Stochastic Programming[M]. 2nd ed. New York: Springer, 2011.
\bibitem{kleywegt} Kleywegt A J, Shapiro A, Homem-de-Mello T. The sample average approximation method for stochastic discrete optimization[J]. SIAM Journal on Optimization, 2002, 12(2): 479--502.
\bibitem{bertsekas} Bertsekas D P. Dynamic Programming and Optimal Control: Vol.~I[M]. 4th ed. Belmont: Athena Scientific, 2017.
\bibitem{rockafellar} Rockafellar R T. Convex Analysis[M]. Princeton: Princeton University Press, 1970.
\bibitem{parisio} Parisio A, Rikos E, Glielmo L. A model predictive control approach to microgrid operation optimization[J]. IEEE Transactions on Control Systems Technology, 2014, 22(5): 1813--1827.
\bibitem{luo} Luo X, Wang J, Dooner M, et al. Overview of current development in electrical energy storage technologies and the application potential in power system operation[J]. Applied Energy, 2015, 137: 511--536.
\bibitem{hejazi} Akhavan-Hejazi H, Mohsenian-Rad H. Optimal operation of independent storage systems in energy and reserve markets with high wind penetration[J]. IEEE Transactions on Smart Grid, 2014, 5(2): 1088--1097.
\bibitem{hong} Hong T, Fan S. Probabilistic electric load forecasting: A tutorial review[J]. International Journal of Forecasting, 2016, 32(3): 914--938.
\bibitem{antonanzas} Antonanzas J, Osorio N, Escobar R, et al. Review of photovoltaic power forecasting[J]. Solar Energy, 2016, 136: 78--111.
\bibitem{weron} Weron R. Electricity price forecasting: A review of the state-of-the-art with a look into the future[J]. International Journal of Forecasting, 2014, 30(4): 1030--1081.
\bibitem{scipy} Virtanen P, Gommers R, Oliphant T E, et al. SciPy 1.0: Fundamental algorithms for scientific computing in Python[J]. Nature Methods, 2020, 17(3): 261--272.
\end{thebibliography}

\newpage
\appendix
\renewcommand{\thesubsection}{\thesection.\arabic{subsection}}
\begingroup
\setstretch{1.22}
\small
% 九、灵敏度分析与组件比较
\section{灵敏度分析与组件比较}\label{sec:sensitivity}
本章回答两个问题：主方法对超参数是否敏感（\S\ref{sec:sens-param}）；各项组件变化怎样影响费用（\S\ref{sec:ablation}）。全部实验与主结果同一 walk-forward 口径、同一执行器、连续执行全年。

\subsection{参数灵敏度}\label{sec:sens-param}
\textbf{场景数 $M$}。固定其余配置（48h+$V(e)$ 主管线），$M\in\{6,12,20,30\}$ 的全年费用为 1375.07、1358.54、1350.75、1352.96 万元：从 6 到 20 条单调改善约 24 万元（更多场景更完整地刻画残差分布，保留水平不再被少数极端路径左右），20 之后进入 $\pm3$ 万元的平台期。$M=30$ 与 $M=20$ 差 0.16\%，说明主结果不依赖场景数的精细选择；取 $M=30$ 兼顾平台稳定与求解时间。

\textbf{前瞻时域}。$M=12$ 下 24h 线性续存与 48h+$V(e)$ 的全年费用为 1362.07 与 1358.54 万元，48h 组件净省 3.53 万元；该增益在 $M=30$ 下为 3.68 万元（表~\ref{tab:components}），两个场景数下方向与量级一致。机理：线性续存费率把午夜库存按固定单价 $v$ 折现，而 $V(e)$ 是凸函数——库存越高边际价值越低；日末光伏富余大的日子，线性口径会高估囤电价值、诱导多买，$V(e)$ 修正了这一偏差（24h 与 48h 配置的平均 24:00 库存分别为 9245.64\,kWh、7200.17\,kWh，不同末端定价对应不同的库存结转策略）。

\textbf{历史窗参数}。续存费率 $v=0.481548$ 并非自由参数，而是“次日 0--5 时均价 $0.4334/\eta$”的推导量；组合权重 $w=0.388815$ 由一月配对误差按最小二乘因果估计（式\eqref{eq:q3combine}），二者都有明确来源且不使用 2--12 月信息。消融配置的 $K_L=6$、$K_P=7$ 与电价窗 $K_p=35$ 由一月实际费用网格搜索选定——全部参数选择只使用一月数据，评分期（2--12 月）从未参与任何调参，杜绝信息泄漏。

\begin{figure}[!htbp]\centering
\safeincludegraphics[width=0.96\textwidth]{fig_f13_sensitivity}
\caption{问题二场景数与前瞻时域的费用比较}\label{fig:fig_f13_sensitivity}
\end{figure}

\subsection{组件比较}\label{sec:ablation}
各项变体均与对应分支的 24h 基线比较，费用增量为“变体费用减基线费用”，见图~\ref{fig:fig_f18_components} 与表~\ref{tab:components}。其中既有启用组件，也有替换预测或场景方法，不把这些实验描述为从 48h 主方法逐项移除组件。负值表示变体更省费，正值表示变体更贵。逐组件解读：

\textbf{48h+$V(e)$（四分支 $-2.4$ 至 $-4.1$ 万元）}：在全部分支都带来节省的时域组件，理由见上文时域段；作为整体组件进入主方法。

\textbf{分解式负载预测（对照“旧均值”，四分支 $+14.5$ 至 $+21.2$ 万元）}：贡献最大的单一组件。日类型均值预测抹平了“周五、周六低负载”的水平切换，而分解式预测用式\eqref{eq:q2beta} 的中位数比率显式建模切换，预测残差方差更小、场景带更窄、保留水平更准。

\textbf{全历史残差池（对照“同类型池”，四分支 $+4.2$ 至 $+5.8$ 万元）}：只用同类型日残差会使早期可用路径减半，场景数不足的代价超过了类型匹配带来的分布贴合收益；残差（预测误差）本身的类型依赖性远弱于负载水平的类型依赖性——切换信息已被点预测吸收。

\textbf{候选计划重评（问 2 增量 $0.00$、问 3 $-0.20$ 万元，问 4-3 为 $-34.96$ 元）}：在 $M=30$ 的这组 24h 对照中，启用候选重评带来的费用变化很小，问 2 与问 4-2 的账单差为零；其收益不能与 48h 组件简单相加，也不计为当前 48h 主结果的费用来源。

\textbf{新价格预测（问 4 对照“旧价格”，4-2 $+0.55$、4-3 $+0.38$ 万元）}：增益存在但有限，与问题四的观察一致——形状易预测导致简单水平规则已能拿到大部分价值。

以上为本评分期上的回顾性组件比较：它支持在该分支保留相应组件，不构成组件收益跨年份可加或普适的证明。


\begin{figure}[H]\centering
\safeincludegraphics[width=0.90\textwidth]{fig_f18_components}
\caption{组件变化的费用增量（相对各分支 24h 基线）}\label{fig:fig_f18_components}
\end{figure}

\begin{table}[H]\centering
\caption{组件变化相对对应单日时域配置的费用增量（元）}\label{tab:components}
\small\renewcommand{\arraystretch}{1.12}\setlength{\tabcolsep}{4pt}
\begin{tabularx}{\textwidth}{@{}>{\hsize=1.5\hsize\linewidth=\hsize}L>{\hsize=.875\hsize\linewidth=\hsize}R>{\hsize=.875\hsize\linewidth=\hsize}R>{\hsize=.875\hsize\linewidth=\hsize}R>{\hsize=.875\hsize\linewidth=\hsize}R@{}}
\toprule
\rowcolor{cumcmgray} 组件变化 & \shortstack{问 2\\单日 13566396.26} & \shortstack{问 3\\单日 13024596.85} & \shortstack{问 4-2\\单日 14270606.71} & \shortstack{问 4-3\\单日 13716972.04} \\
\midrule
启用 48h 与次日价值 & -36761.01 & -35046.21 & -40706.98 & -23960.45 \\
启用候选计划重评 & 0.00 & -1974.51 & 0.00 & -34.96 \\
改用同类型日场景池 & 42116.91 & 47354.21 & 57738.37 & 53857.16 \\
改用旧均值负载预测 & 204242.11 & 145252.73 & 211566.29 & 144922.98 \\
改用旧价格预测 & --- & --- & 5524.96 & 3816.14 \\
\bottomrule
\end{tabularx}
\par\vspace{3pt}{\footnotesize 费用增量为变体减单日时域配置，负值表示节省；”---”表示该分支不适用。}
\end{table}


\FloatBarrier
\endgroup
\section{支撑材料文件列表与源程序}
支撑材料包含：全部源程序（\texttt{src/} 目录：数据读取、优化核、预测与场景、DP 价值执行器、结算与结果导出）、结果文件 result1--result4-3.xlsx、全年实验汇总与独立复核记录、AI 工具使用详情.pdf。核心求解环境为 Python 3.13 + SciPy（HiGHS 线性规划）。以下给出优化核与 DP 价值执行器的关键源代码，其余模块见支撑材料。

\lstset{language=Python,basicstyle=\ttfamily\scriptsize,breaklines=true,
  numbers=left,numberstyle=\tiny,frame=single,tabsize=4,
  showstringspaces=false,columns=flexible}

\subsection{预报时刻组合的精确结果（单日时域配置）}
\begin{table}[!htbp]
\centering\small
\caption{问3发布时刻组合与相邻组合的费用变化}\label{tab:issues}
\begin{tabular}{lrrr}
\toprule
发布时刻/h & 总费用/元 & 紧急费/元 & 相邻节省/元 \\
\midrule
0 & 13351182.68 & 579245.64 & --- \\
0,6 & 13312853.41 & 353878.30 & 38329.27 \\
0,6,12 & 13075459.34 & 265785.37 & 237394.07 \\
0,6,12,18 & 13024596.85 & 192075.49 & 50862.49 \\
\bottomrule
\end{tabular}
\end{table}

\FloatBarrier
\subsection{优化核（optimizer.py，节选：确定性 LP 与 SAA）}
\lstinputlisting[firstline=204,lastline=290]{../../src/optimizer.py}

\subsection{凸折线价值函数与动态规划（value\_dp.py，节选）}
\lstinputlisting[firstline=328,lastline=428]{../../src/value_dp.py}

\subsection{DP 价值执行器（executor.py，节选）}
\lstinputlisting[firstline=211,lastline=234]{../../src/executor.py}
\end{document}
